\chapter{Ergebnisse}

Im Ergebniskapitel werden die Resultate der Arbeit präsentiert — sachlich, objektiv und ohne Wertung. Die Interpretation und kritische Einordnung erfolgt erst im nächsten Kapitel (Diskussion). Diese Trennung ist wichtig, damit der Leser die Fakten zunächst unvoreingenommen aufnehmen kann.

\section{Auswertung}

Typische Inhalte sind: Messergebnisse, Simulationsdaten, Performance-Metriken einer Software oder die Ergebnisse einer Nutzerstudie. Grafiken und Tabellen sind hier besonders wichtig, um Daten übersichtlich zu präsentieren.

Achte darauf, dass jede Abbildung und Tabelle im Text referenziert und kurz beschrieben wird. Nicht nur zeigen, sondern auch beschreiben, was zu sehen ist — der Leser soll nicht raten müssen, worauf er achten soll.

\section{Zusammenfassung der Ergebnisse}

Am Ende des Kapitels kann eine kurze Zusammenfassung stehen, die die wichtigsten Resultate auf einen Blick bündelt. Das erleichtert dem Leser den Übergang zur Diskussion.